\scanpage{44}\deleted{\unsure{48|49}}\mspage{50}\deleted{\unsure{a|2}}
\deleted{Je vais généraliser légèrement le \deleted{\unsure{lemme}} \inserted{Cor. 1}}

\uline{Cor. 5} (g) \deleted{Soit} Soit $f:X\to Y$ pleinement fidèle dans Diag, et $F\in\mathbb D(X)$. Alors l'hom d'adjonction
\[
 f^*f_*(F)\longrightarrow F
\]
est un \uline{iso}.

En effet, il suffit de le vérifier fibre par fibre (Der 2). Soit, $x\in X$, \inserted{$y=f(x)\in Y$,} alors
\[
 f^*f_*(F)_x\simeq f_*(F)_y\underset{\text{Dér 5}}{\simeq}\deleted{\unsure{X_{/x}}}\Gamma_{\mathbb D}(X_{/y},F)
\]
or \deleted{\unsure{$X_{/y}$}} à cause de la pleine fidélité $X_{/y}\simeq X_{/x}$. \deleted{\unsure{Mais} Mais \unsure{j'ai|il} \unsure{oublié} l'\unsure{axiome}}

% slanted
\margin{Va résulter de \uline{Der 6} (g), plus fort ci-dessous}
% slanted
\margin{Axiome : \deleted{\unsure{limites}}\deleted{\unsure{coefficients}} la cat. d'indices à objet final (coeff. \uline{quelc.})}
\deleted{\unsure{Dér} \unsure{4}$_0$ (g)} : \inserted{\uline{Cor. 3} (\unsure{remonte|renvoi à} p. 48)} Soit $Z\in\mathrm{Diag}$, supposons que $Z$ a un objet final, \inserted{$e_Z$,} \deleted{plus généralement que le foncteur $\underleftarrow{\lim}_{Z}$ \unsure{ensembliste} (pour les préfaisceaux d'ens sur $Z$) transforme épi en épi (exactitude à droite de $\underleftarrow{\lim}$). Alors \unsure{Alors}} Soit $p:e\to Z$ le foncteur correspondant. Alors \deleted{\unsure{le}} l'hom canonique (pour $F$ dans $\mathbb D(Z)$)
\diagram{\begin{tikzcd}[ampersand replacement=\&,row sep=large]
 Z\arrow[d,"f"]\\ e\arrow[u,bend left=24,dashed,"p"]
\end{tikzcd}}
\[
 f_*(F)\to p^*(F)
\]
est un iso.

Cet hom \deleted{se déduit} est \unsure{le particulier|un particulier} de celui-ci, pour
% slanted
\margin{déjà envisagé dans Cor. 2 p. 47}
\diagram{\begin{tikzcd}[ampersand replacement=\&,row sep=small]
 Z'\arrow[rr,"p"]\arrow[dr,"f'"'] \& \& Z\arrow[dl,"f"]\\
 \& S \&
\end{tikzcd}}
(en faisant $Z'=S=e$, $f'=\mathrm{id}_e$)
\[
 f_*(F)\to f'_*(p^*(F))\qquad(F\in\mathrm{Ob}\,\mathbb D(Z))
\]
et c'est un cas particulier du morphisme de
% archive sheet number: 43
